\documentclass[11pt]{article}
\usepackage[margin=1.1in]{geometry}
\usepackage{xcolor}
\usepackage{hyperref}
\hypersetup{colorlinks=true, urlcolor=blue, linkcolor=black}
\usepackage{enumitem}
\setlength{\parindent}{0pt}
\setlength{\parskip}{6pt}

\newcommand{\rc}[1]{\medskip\noindent\textbf{\textcolor[rgb]{0.35,0.35,0.35}{#1}}\par\nopagebreak}
\newcommand{\reply}{\medskip\noindent\textbf{Response.}\ }
\newcommand{\where}[1]{\par\smallskip\noindent\emph{Where in the manuscript:} #1\par}

\title{\vspace{-2.2em}Response to Reviewers\\[2pt]
\large IoT-66559-2026 (R1)\\[4pt]
\normalsize submitted as: FedKAN-IDS: Heterogeneity-Robust Federated Kolmogorov--Arnold
Networks for NetFlow-Based IoT Intrusion Detection\\[2pt]
\normalsize revised to: Tuning Confounds the Comparison: What a Controlled Study of
Federated KANs for NetFlow Intrusion Detection Actually Shows}
\author{}
\date{}

\begin{document}
\maketitle
\vspace{-3em}

We thank both reviewers for reports that were specific enough to act on. Two of the
requested experiments changed our conclusions rather than confirming them, and we
report that plainly below: Reviewer 2's comment 7 in particular led us to a result
that weakens the manuscript's original headline claim. We have rewritten the paper
around what the new measurements actually support.

\textbf{Summary of what changed.}
\begin{itemize}[leftmargin=1.4em, itemsep=2pt]
  \item Every architecture is now tuned over a learning-rate grid before being
        compared ($1{,}200$ runs for the sweep, $1{,}500$ in total), and all headline numbers are re-reported under that
        protocol. The title, abstract and contributions changed as a result: the accuracy
        and variance claims of the first version are withdrawn, and the paper now reports a
        robustness-versus-compute trade-off.
  \item All experiments were re-run on a single machine so that no number in the
        paper is compared across hardware.
  \item Accuracy is now summarised by the mean over the final five rounds rather
        than the last round alone.
  \item Precision, recall, MCC and false-positive rate are reported for every cell.
  \item Communication cost is reported as bytes-to-target and rounds-to-target, not
        only total bytes after a fixed round budget.
\end{itemize}

\hrule
\section*{Reviewer 1}

\rc{R1.1 Novelty clarification: ``The contribution is primarily empirical rather than
methodological'': articulate scientific insights beyond empirical KAN vs.\ MLP
comparisons.}
\reply 
We agree that ``we compared two architectures'' is not by itself a scientific insight, and
in revising we found that the insight the study actually supports is a different and
sharper one than the submitted version claimed. Our learning-rate sweep (see R2.7) shows
that the accuracy and variance advantages attributed to KANs in the FedKAN literature,
including in our own submitted manuscript, are confounded with \emph{how much
hyper-parameter tuning each architecture received}. Once each architecture is given its
own best learning rate, the mean advantage and the variance advantage both largely
disappear on the cell that produced our headline number. What survives is a distinct and
practically relevant property: the KAN reaches within a fraction of a point of its best
accuracy across a wide band of learning rates, whereas the parameter-matched MLP requires
tuning to be competitive at all. For federated deployment across IoT gateways, where
per-client hyper-parameter search is not available, insensitivity to the learning rate is
arguably the more useful property, and it is a claim about \emph{optimisation behaviour}
rather than about representational capacity. The revised paper is built around that claim,
and the title and abstract have changed accordingly.
\where{}

\rc{R1.2 Expand baselines: add comparisons with FedProx, SCAFFOLD, MOON, FedNova, or
FedDyn beyond FedAvg-based MLPs.}
\reply Added, on the cell that produced our headline number, ten seeds each. FedProx is the one that
matters and it works against our original claim: it lifts the parameter-matched MLP from
$88.72\%$ to $94.92\%$, a gain of $6.19$~pp, and cuts its seed variance from $14.21$ to
$2.67$~pp, while lifting FedKAN-8 by only $0.44$~pp. Under FedProx the two architectures are
level and the sign in fact reverses: $-0.26$~pp in the KAN's disfavour ($p=0.82$). Our
submitted comparison therefore under-served the baseline in two independent ways, a shared
learning rate and a FedAvg-only aggregator, and correcting either one removes most of the gap.
MOON gives no benefit to either architecture ($91.17\%$ and $91.14\%$, $p=0.32$ and $p=0.55$).

We should add that this table was originally measured at exactly the shared step size the rest of
the revision argues against, which would have been a fair objection. We therefore re-ran FedProx
with each architecture at its own best rate: FedKAN-8 reaches $93.85\%$ and the MLP $95.72\%$, a
difference of $-1.87$~pp. The two corrections compound rather than cancel, and the revised
Table~\ref{tab:compound} reports all three configurations side by side.

We must be candid about SCAFFOLD. Our implementation does not converge here: $49.80\%$ for
FedKAN-8, which is chance on a binary task, and $73.32\%$ for the MLP. SCAFFOLD's control
variates are derived for local SGD, and every experiment in this paper uses Adam; we believe
that mismatch is the cause but we have not proved it. We therefore report the runs and
explicitly draw \emph{no} conclusion about SCAFFOLD from them, rather than present a number we
suspect is our own error as evidence about the method. FedNova and FedDyn were not run; we
prefer to say so than to add them hastily.
\where{}

\rc{R1.3 IDS metrics: include Precision, Recall, F1-score, MCC, ROC-AUC, or False
Positive Rate for comprehensive evaluation.}
\reply Precision, recall and per-class F1 were already computed for every run; we
have added them to the results tables together with MCC and the false-positive rate,
which we derive exactly from the stored per-class recalls and accuracies. We verify
the derivation rather than assert it: reconstructing each confusion matrix and
recomputing precision reproduces the stored value for $163/163$ binary runs.
We do \emph{not} report ROC-AUC, and we prefer to say so explicitly rather than pass
over the request: our runs store predicted labels, not continuous scores, so ROC-AUC
cannot be recovered from the released artefacts.

\where{}

\rc{R1.4 Communication costs: beyond parameter transmission, add transmitted bytes,
communication rounds, convergence efficiency, or training time.}
\reply We now report bytes and rounds \emph{to reach a fixed accuracy target} rather than total bytes
after a fixed round budget, because the latter assumes both architectures must run the full
schedule. At a target of $90\%$ of the best accuracy any architecture reaches, FedKAN-8 needs
$1.32$~MB of uplink against $0.10$~MB for FedAvg-MLP-32; at $94\%$ the figures are $1.04$ and
$0.54$~MB. We also report the number of seeds that reach the target at all, which is the
column that favours the KAN: $9/10$ and $8/10$ against $7/10$ and $6/10$. Training time and
per-round wall-clock are included, and inference cost is answered under R1.5.
\where{}

\rc{R1.5 Deployment claims: support Raspberry Pi statements with experimental evidence
or present them more cautiously as potential future work.}
\reply We have removed the Raspberry Pi statements rather than soften them, because we do not have
that hardware and a number we cannot measure should not appear. In their place we report what
we can measure: single-thread CPU inference with batch size $1$, which is the closest
deployment-like regime available to us. At matched parameters ($3{,}280$ against $3{,}362$)
FedKAN-8 costs $0.0995$~ms per sample against $0.0050$~ms for the parameter-matched MLP, a
factor of $19.9$, and $1.690$~ms against $0.045$~ms for a batch of $1000$, a factor of $37.9$.
So that a reader can convert these to their own gateway we also report an anchor on the same
machine: a $256\times256$ matrix multiply takes $0.024$~ms. Parameter parity is not cost
parity, and this now appears in the abstract.
\where{}

\rc{R1.6 Generalization: present class-distribution skew and FedKAN performance as an
empirical observation rather than a universally established property.}
\reply 
We have removed the generalising language. Every statement linking class-distribution skew
to FedKAN performance is now scoped explicitly to the three NetFlow-v2 datasets and the
partitioning regimes we ran, and is presented as an empirical observation on that sample
rather than an established property. We have also removed the sentences in the
introduction and conclusion that projected the pattern onto unseen datasets. Reviewer~1 is
right that three datasets cannot establish a monotone relationship, and our revised
analysis (R2.7) gives a further reason for caution: the direction of the effect is
sensitive to a hyper-parameter choice we had held fixed.
\where{}

\rc{R1.7 Presentation: increase font sizes, improve confidence interval visibility,
simplify long sentences.}
\reply 
All figures have been redrawn: axis and tick label sizes are increased, confidence
intervals are drawn as shaded bands with explicit whisker caps at the mean, and the
colour palette is now colour-blind safe and locked per method across every figure so that
one architecture keeps one colour throughout. We have also gone through the manuscript for
overlong sentences. We note one specific defect the reviewer's comment led us to: a
sentence in the conclusion had lost its punctuation in the submitted PDF and read as
``a reproducible failure mode seed~17's Dirichlet partition drives\ldots''. That sentence
has been repaired.
\where{}

\hrule
\section*{Reviewer 2}

\rc{R2.1 Theoretical assumptions: Theorem 1 assumes $L$-smoothness and $\mu$-strong
convexity, which cross-entropy loss with deep networks does not satisfy globally.}
\reply 
The reviewer is correct, and on re-examining the section we found a second and more
serious problem that we would rather state ourselves than have discovered later.

First, on the assumptions as raised: cross-entropy loss composed with either a KAN or an
MLP is not globally $L$-smooth or $\mu$-strongly convex, so Theorem~\ref{thm:conv} does
not describe our training runs globally. We now state this in the theorem's preamble, and
we have relabelled the result as describing behaviour in a neighbourhood of a local
minimum under those assumptions, which is the standard reading of the FedAvg bound we
adapt~\cite{li2020convergence}.

Second, and more importantly: Corollary~2 in the submitted manuscript does not follow from
Theorem~\ref{thm:conv}. The bound's heterogeneity term $\frac{2L}{\mu^2 T}(L+\Gamma_\alpha)$
applies identically to both architectures, and the spline-approximation residual
$\mathcal{O}(G^{-(k+1)})$ is an \emph{additional} error term for the KAN, not a mechanism
that shields it from $\Gamma_\alpha$. The partial derivative with respect to
$\Gamma_\alpha$ is the same for both. The corollary therefore asserted the robustness
conclusion rather than deriving it. Since our revised experiments no longer support that
conclusion empirically either (R2.7), we have removed Corollary~2 rather than repair it,
and the convergence analysis is now presented only as what it establishes: that the
spline bias enters additively and does not scale with heterogeneity.
\where{}

\rc{R2.2 Missing baselines: add communication-efficient FL methods (FedPAQ, Top-$k$
SGD, or CG-FKAN); at minimum one compression baseline.}
\reply Added FedPAQ with top-$k$ sparsification at $k=10\%$, with byte accounting that charges four
bytes for each retained value and four for its index. At that budget the uplink falls from
$7.69$ to $1.64$~MB for FedKAN-8, a factor of $4.7$, but accuracy falls to $70.59\%$
($p=0.003$) and for the MLP to $55.52\%$ ($p<0.001$). We report this as a negative result at
this compression level rather than tune $k$ until it flatters us: under Dirichlet
$\alpha=0.1$, where client updates already disagree, discarding $90\%$ of each update
compounds the disagreement. CG-FKAN is discussed in related work but not reproduced, since it
requires a grid-sparsification schedule we could not reconstruct from the paper; we state that
rather than approximate it.

We also note the accounting point, because it is easy to miss: SCAFFOLD doubles the uplink
($15.38$ against $7.69$~MB) because each client returns a control variate as well as a model
delta. A comparison that charges every algorithm one model per round would make SCAFFOLD look
free.
\where{}

\rc{R2.3 Seed-17 pathology: provide class and feature distributions per client, why
F-KAN-16 escapes the trap and MLP does not, and whether other seeds behave similarly.}
\reply We measured the partition statistics the reviewer asks for, on all ten seeds so that seed~17
can be ranked rather than merely described, and the result is that \emph{no statistic we
measured singles it out}. On NF-ToN-IoT-v2 seed~17 ranks second of ten by client-size Gini
($0.7235$, behind seed~31 at $0.7600$, which is not pathological), ninth of ten by mean
per-client label entropy (seed~29 is lower still and performs normally), and its
class-proportion spread of $1.0$ is shared by eight of the ten seeds. Rank correlations against
final accuracy are weak and not significant: Spearman $+0.42$ ($p=0.23$) for entropy, $-0.30$
($p=0.40$) for Gini, $+0.61$ ($p=0.06$) for smallest client size. We therefore report the
failure mode as reproducible but unexplained rather than offer a mechanism the data does not
support, and we have removed the sentence attributing it to an interaction between the
partition and the class distribution.
\where{}

\rc{R2.4 Hardware profiling: include preliminary measurements on Raspberry Pi or
similar hardware; RTX 4090 times are not representative of IoT deployment.}
\reply Answered under R1.5: single-thread CPU profiling with a conversion anchor, in place of the
RTX~4090 numbers, which we agree answer a different question (centralised training throughput)
and are labelled as such. We did not fabricate Raspberry Pi figures.
\where{}

\rc{R2.5 Sensitivity analysis: vary grid size $G \in \{3,5,10\}$, spline order
$k \in \{3,4,5\}$, and hidden width beyond 8 and 16.}
\reply Added, on ten seeds each, with parameter counts alongside so that a capacity effect cannot be
mistaken for a spline effect. Spline order is irrelevant here: $k{=}4$ gives $+0.10$~pp and
$k{=}5$ gives $-0.12$~pp against $k{=}3$. Grid size matters and we re-ran it on ten seeds rather than the single run of the
first version: $G{=}3$ costs $0.87$~pp and $G{=}10$ gains $1.93$~pp on NF-BoT-IoT-v2
($p=0.034$) and $3.89$~pp on NF-ToN-IoT-v2 ($p=0.095$), with the false-positive rate improving
from $8.36$ to $4.43\%$ and from $21.29$ to $16.01\%$. We should correct something here: an
intermediate draft of this response cited a $84\%$ false-positive rate for $G{=}10$, which came
from the single run available at the time and is not representative. With ten seeds $G{=}10$ is
better on every metric we report. It is still not free, adding $50\%$ parameters, which breaks
the parameter parity the comparison rests on, and roughly ten times the training time, and it
does not repair the seed-17 partition, which moves from $51.79$ to $56.12\%$. Hidden width is
non-monotone and this surprised us: $h{=}4$, at half the parameters of our default, is
$0.98$~pp \emph{better} than $h{=}8$, while $h{=}16$ and $h{=}32$ give $+0.80$ and $+0.63$~pp
at two and four times the parameters. Within this family more capacity does not help, which is
consistent with the regularisation reading of the spline basis and inconsistent with a
capacity reading.
\where{}

\rc{R2.6 Local convexity: provide empirical evidence (e.g.\ Hessian spectrum analysis)
supporting the local convexity assumption.}
\reply Measured, and the answer is unfavourable to the theorem. We ran Lanczos with full
re-orthogonalisation on the exact Hessian-vector product at the converged federated solution
(NF-BoT-IoT-v2, binary, $50$ rounds, $8192$ held-out samples, $40$ Ritz directions). FedKAN-8
gives $\lambda_{\max}=88.5$, $\lambda_{\min}=-0.134$, with $45\%$ of Ritz values negative;
the parameter-matched MLP gives $\lambda_{\max}=0.52$, $\lambda_{\min}=-1.8\times10^{-3}$,
with $27.5\%$ negative. Negative curvature at the solution means the iterates settle at a
saddle region, so strong convexity fails not only globally, as the reviewer notes, but locally
as well. We state this in the revised theory section, and it is part of why we removed
Corollary~2. We also note what the measurement does \emph{not} show: the two condition numbers
are the same order ($662$ against $285$), so the Hessian spectrum does not explain any
performance difference between the architectures.
\where{}

\rc{R2.7 Hyperparameter justification: clarify whether Adam with $\eta = 10^{-2}$ was
tuned per architecture, and whether KANs require different learning rates.}
\reply This comment was correct and acting on it changed the paper's central claim. The submitted
study used $\eta=10^{-2}$ for every architecture. We swept six values spanning two and a half
orders of magnitude, for all four architectures, on all ten seeds, on all four cells:
$1{,}200$ runs on one machine.

The parameter-matched MLP's best step size on NF-BoT-IoT-v2 is $10^{-1}$, ten times the shared
value. Giving it that value gains it $5.39$~pp; the same freedom gains FedKAN-8 $0.65$~pp. The
headline gap falls from $+5.49$ to $+0.76$~pp ($p=0.74$ paired, bootstrap CI
$[-3.39,+4.99]$). Across the four cells the advantage ranges from $-0.49$ to $+6.23$~pp with
no consistent sign. The variance-reduction ratio behaves the same way: $2.53\times$ becomes
$1.05\times$ once both sides are tuned, and is already $0.71\times$, $0.33\times$ and
$0.87\times$ on the other cells.

Yes, KANs need a different learning rate from MLPs, and that turns out to be the paper's real
finding rather than a nuisance: KAN is $2.7$ to $5.4$ times less sensitive to the choice.
Section~\ref{sec:lrsweep} is new, the title and abstract are rewritten around this, and the
accuracy and variance claims of the first version are withdrawn.
\where{}

\rc{R2.8 Seed-17 characterization: detail the client data distributions that make this
partition uniquely challenging for KAN-8.}
\reply Answered together with R2.3 above: the per-client distributions are now reported for all ten
partitions, and they do not identify seed~17.
\where{}

\end{document}
